% !TeX root = ../main.tex
% ---------------------------------------------------------------------------
% Shotcrete and Sprayed Concrete Linings
% Scope extension: expands the shotcrete section of Chapter 10 ( Chapter~\ref{ch:construction}).
% Key sources: ACI PRC-506; EN 14487; ASTM C1116; Austrian Society for
% Construction Technology SCL guideline; Thomas, Sprayed Concrete Lined Tunnels.
% ---------------------------------------------------------------------------
\chapter{Shotcrete and Sprayed Concrete Linings}
\label{ch:shotcrete}


\section{Role and spraying process}
\label{sec:role-and-spraying-process}

\subsection{The role of sprayed concrete}
\label{sec:shotcrete-role}
\scopenote{From surface sealing to full permanent structural lining. How
sprayed concrete came to carry primary, and increasingly permanent, structural
function.}

\subsection{Wet mix and dry mix processes}
\label{sec:wet-dry-mix}
\scopenote{Process differences, output rates, rebound, dust, consistency and
the near-universal adoption of wet mix on major projects.}

\section{Mix design and reinforcement}
\label{sec:mix-design-and-reinforcement}

\subsection{Mix design and admixtures}
\label{sec:shotcrete-mix}
\scopenote{Binder content, aggregate grading, water-binder ratio, plasticisers
and hydration control agents; alkali-free accelerators and their effect on
long-term strength.}

\subsection{Fibre reinforcement}
\label{sec:fibres}
\scopenote{Steel and macro-synthetic fibres, dosage, orientation effects,
residual flexural strength classes, and energy absorption panel testing.}

\section{Early-age and structural behaviour}
\label{sec:early-age-and-structural-behaviour}

\subsection{Early-age strength development}
\label{sec:early-age}
\scopenote{The J1--J3 early strength classes, penetrometer and stud driving
tests, and why the early-age curve, not the 28-day strength, governs
excavation cycle time.}

\subsection{Structural behaviour of a sprayed lining}
\label{sec:scl-behaviour}
\scopenote{The lining as a thrust-dominated shell; thrust--moment interaction
diagrams; ageing stiffness and creep; and how load is shared with the ground
over time.}

\section{Lining design}
\label{sec:lining-design}

\subsection{Design of primary and permanent linings}
\label{sec:scl-design}
\scopenote{Single-shell versus double-shell philosophy, composite and
sacrificial primary linings, and the verification required before a sprayed
lining is accepted as permanent works.}

\subsection{Yielding elements for squeezing ground}
\label{sec:yielding-elements}
\scopenote{Deformation slots and compressible elements that let the lining
converge without overstress; the design principle behind the ground reaction
curve of Chapter~\ref{ch:ccm} being met at lower pressure.}

\subsection{Reinforced ribs of sprayed concrete}
\label{sec:rrs}
\scopenote{The NMT RRS system, rib geometry and spacing, and its position
between plain shotcrete and steel sets.}

\section{Steel elements and substrate}
\label{sec:steel-elements-and-substrate}

\subsection{Lattice girders and steel sets}
\label{sec:girders}
\scopenote{Function as profile control and early support, embedment in
shotcrete, and the debate over their real structural contribution.}

\subsection{Substrate preparation, bond and layering}
\label{sec:shotcrete-bond}
\scopenote{Surface cleaning, bond strength, cold joints between layers, and
drainage behind the lining.}

\section{Durability and quality control}
\label{sec:durability-and-quality-control}

\subsection{Durability and fire performance}
\label{sec:shotcrete-durability}
\scopenote{Permeability, cover to reinforcement, aggressive groundwater
exposure, and explosive spalling under fire with polypropylene fibre
mitigation.}

\subsection{Spraying practice, rebound and quality control}
\label{sec:spraying-practice}
\scopenote{Nozzle distance and angle, robotic versus hand spraying, rebound
quantification, thickness control, coring and panel testing.}
